\section{Perspectiva teórica de Meta-RL}

\subsection{Meta-RL como POMDP}

\begin{importantbox}{Meta-RL vista como POMDP}
Meta-RL puede entenderse como hacer RL dentro de un POMDP más amplio:
\begin{itemize}
    \item El estado oculto incluye la identidad de la tarea
    \item La observación es el estado actual junto con la experiencia histórica
    \item La política óptima necesita inferir la identidad de la tarea durante la interacción (inferencia bayesiana)
\end{itemize}
Desde esta perspectiva, los métodos context-based en esencia aproximan la política óptima de este POMDP.
\end{importantbox}

\subsection{Muestreo posterior y Thompson Sampling}

En el caso ideal, el comportamiento óptimo de meta-RL debería parecerse a Thompson Sampling:
\begin{enumerate}
    \item Mantener una creencia posterior sobre la identidad de la tarea $p(\mathcal{T} | \text{history})$
    \item Muestrear de la distribución posterior una hipótesis de tarea
    \item Actuar según la política óptima bajo esa hipótesis
    \item Actualizar la distribución posterior después de recolectar datos nuevos
\end{enumerate}

\subsection{Resumen de la sección}

Desde el punto de vista teórico, Meta-RL puede entenderse de manera unificada como la toma de decisiones óptimas en un POMDP con incertidumbre sobre la tarea, lo cual ofrece un marco basado en principios para entender y mejorar los métodos de meta-RL.

